\section{模型求解与结果}
% 写作提示：按子问题交代数据处理、算法步骤与参数依据、计算环境、关键结果及其单位，并解释结果如何回答题目。
% 每张图表都应在正文引用和解读；下方回归、随机森林、LASSO、遗传算法及全部数值仅为示例，不得直接沿用。
\subsection{数据预处理}
采用Z-score标准化处理：$x' = \frac{x - \mu}{\sigma}$。

\subsection{模型建立}
采用多元线性回归作为基础模型：
\begin{equation}
y = \theta_0 + \theta_1 x_1 + \theta_2 x_2 + \cdots + \theta_n x_n + \varepsilon
\end{equation}
同时引入随机森林模型进行对比分析，通过集成多棵决策树提高精度和鲁棒性。

\subsection{求解结果}
随机森林的测试集$R^2$达到0.896，优于线性回归的0.834。
\subsection{特征选择}
采用LASSO回归筛选关键特征：
\begin{equation}
\min_{\theta} \frac{1}{2N} \sum_{i=1}^{N} (y_i - \theta^T x_i)^2 + \lambda \|\theta\|_1
\end{equation}

\subsection{模型优化}
基于筛选的关键特征建立优化预测模型，引入交叉项捕捉交互效应。
优化后测试集$R^2$从0.896提升至0.923，RMSE降低约15\%。
\subsection{优化模型}
将问题建模为约束优化问题：
\begin{align}
\min \quad & f(\mathbf{x}) = \sum_{i=1}^{n} c_i(x_i) \\
\text{s.t.} \quad & g_j(\mathbf{x}) \leq 0, \quad h_k(\mathbf{x}) = 0
\end{align}

\subsection{算法求解}
采用遗传算法（种群规模200，迭代100代）进行求解，算法快速收敛到稳定解。

\subsection{结果}
得到最优策略方案：$x_1=0.75$，$x_2=1.20$，$x_3=0.45$，$x_4=0.88$。
